\chapter{چارچوب گذار: مجلس مهستان}

\begin{modernbox}{یادداشت مهم درباره‌ی اعداد و نسبت‌ها}
تمام اعداد، نسبت‌ها، آستانه‌ها و زمان‌بندی‌های مشخصی که در این سند ذکر می‌شوند (مانند ۳۰۰ نماینده، ۳٪ آستانه و غیره) \textbf{پیشنهادی برای بحث و شروع گفتگو} هستند. منظور اصلی این سند، نشان دادن نیاز به تعریف چنین پارامترهایی است — نه تثبیت خود آن اعداد.
\end{modernbox}

نهاد محوری گذار یک \textbf{مجلس مردم‌محور منتخب} (مجلس مهستان) است که در اسرع وقت پس از آغاز دوره‌ی گذار تشکیل خواهد شد.

\section{فرایند انتخابات و تشکیل مجلس}

\begin{enumerate}
    \item جریان‌ها و ائتلاف‌ها فهرست کاندیداهای خود را معرفی می‌کنند. هر جریانی که فهرست ملی معرفی می‌کند، موظف است پیش‌نویس قانون اساسی موقت پیشنهادی خود را نیز منتشر کند.
    \item ثبت‌نام رسمی در وب‌سایت وزارت کشور و سایر پایگاه‌های اطلاع‌رسانی منعکس خواهد شد.
    \item اجرای انتخابات به‌صورت غیرمتمرکز توسط فرمانداری‌ها و شهرداری‌ها با نظارت ناظران بین‌المللی و جامعه‌ی مدنی خواهد بود.
\end{enumerate}

\section{ترکیب مجلس مهستان}
مجلس مهستان متشکل از \textbf{۳۰۰ نماینده} خواهد بود:
\begin{itemize}
    \item \textbf{۲۴۰ کرسی استانی:} بر اساس جمعیت هر استان، با سیستم نسبی-فهرستی.
    \item \textbf{۴۰ کرسی فهرست ملی:} برای جریان‌هایی که فهرست سراسری دارند.
    \item \textbf{۲۰ کرسی تضمینی اقلیت‌ها:} برای تضمین حضور صدای اقلیت‌های قومی، زبانی و مذهبی.
\end{itemize}

\begin{highlightbox}
\textbf{ضمانت اجرایی سهمیه‌ی جنسیتی (قاعده‌ی ۱ به ۳):}
ثبت‌نام فهرست‌ها منوط به رعایت این شرط است: در میان هر سه نفر متوالی در فهرست، حداقل یک نفر باید از جنسیت کمتر‌نمایندگی‌شده باشد. در غیر این صورت، نفر سوم با بالاترین رأیی که به یک زن داده شده جایگزین می‌شود.
\end{highlightbox}

\section{هیأت رئیسه و ساختار داخلی}
مجلس مهستان بلافاصله کمیسیون‌های تخصصی دائمی زیر را تشکیل خواهد داد:
\begin{itemize}
    \item کمیسیون قانون اساسی موقت
    \item کمیسیون عدالت انتقالی و حقوق بشر
    \item کمیسیون اصلاح بخش امنیتی و نظامی
    \item کمیسیون اقتصاد و بودجه
    \item و سایر کمیسیون‌های تخصصی...
\end{itemize}

\section{تشکیل کابینه‌ی اجرایی موقت}
در صورت عدم تشکیل هیچ ائتلافی، مجلس مهستان برای هر پست کابینه به‌طور جداگانه رأی‌گیری می‌کند و افرادی که بالاترین رأی را کسب کنند، به‌طور خودکار انتخاب می‌شوند. این سازوکار از خلأ اجرایی جلوگیری می‌کند.
